\begin{table}[htbp]
\centering
\caption{Gap Analysis: Current EPC vs.\ Potential--Current Gap}
\label{tab:gap_analysis}
\begin{tabular}{lrrrl}
\toprule
Treatment & Current EPC coef & Gap coef & $|$Ratio$|$ & Interpretation \\
\midrule
For-profit (all) & +0.42*** & -0.07 & 0.16 & Stock selection \\
Tax haven (all) & +1.61*** & -0.63*** & 0.39 & Hybrid \\
Foreign for-profit & +1.64*** & -0.64*** & 0.39 & Hybrid \\
Non-profit & +3.15*** & -2.26*** & 0.72 & Predominantly operational \\
Public sector & +3.35*** & -2.57*** & 0.77 & Predominantly operational \\
\bottomrule
\end{tabular}
\begin{minipage}{0.9\textwidth}
\vspace{4pt}
\footnotesize
\textit{Notes:} PSM with subclass FE, LA geography, Baseline spec, base core.
\textbf{Speculative.} The cross-sectional gap is not a clean test of operational
investment; panel data would be needed to identify retrofit activity directly.
Current EPC coef: coefficient on \texttt{current\_energy\_efficiency}.
Gap coef: coefficient on \texttt{energy\_efficiency\_potential\_gap} (= potential EPC $-$ current EPC).
Ratio = $|$gap coef $/$ current EPC coef$|$.
A ratio near 0 means both current and potential rise together (selection);
a ratio near 1 means the gap shrinks (operational investment).
$^{***}p<0.001$; $^{**}p<0.01$; $^{*}p<0.05$.
Source: \texttt{results\_enriched\_LA.csv}.
\end{minipage}
\end{table}
